\section{Day and night system}
\label{sec:day}
\here{8}
The sun follows a solar-position model for the camp's latitude (Khulna, $\varphi=22.8^\circ$, early October,
declination $\delta=-4.6^\circ$). It rises in the east, crosses the southern sky and sets in the west. Every lighting
parameter is a function of its elevation $e$.
\begin{gather}
\sin e=\sin\varphi\sin\delta+\cos\varphi\cos\delta\cos H,\quad H=15^\circ(\text{hour}-12)\label{eq:sun}\\
T_c(e)=\exp\!\big(-\tau_c\,(m(e)-1)\big),\quad \tau=(0.06,\,0.11,\,0.21)\nonumber
\end{gather}
with the Kasten--Young air mass \cite{kasten1989} $m(e)=1/\big(\sin e+0.50572\,(e+6.07995^\circ)^{-1.6364}\big)$. The sun's key light is
$3.2\,s(e)\,T(e)$ with $s$ a smoothstep from $-1^\circ$ to $10^\circ$; blue has the largest optical depth, so it is
removed first near the horizon and the light turns orange.

\begin{figure}[htbp]
\centering\plot{p_day_cycle}{\linewidth}
\caption{The solar model over 24 hours (Eq.~\eqref{eq:sun}): (a) sun elevation with the twilight bands where the
phases switch, (b) the key light per channel, (c) the zenith, horizon and key-light colours the renderer uses, read from
the renderer's export at 15-minute steps.}
\label{fig:daycycle}
\end{figure}

The phases switch at night below $-6^\circ$ and day above $4^\circ$ (Figure~\ref{fig:daycycle}a). Red outlasts blue at
both ends of the day (b): that is the sunset colour of Figure~\ref{fig:hero}a and the warm horizon band at 6:00 and 18:00 in (c), produced by the formula.
\begin{figure}[htbp]
\centering
\begin{subfigure}{0.325\linewidth}\raw{scene/dawn_withdraw}{\linewidth}\caption{Dawn: the intruders withdraw.}\end{subfigure}\hfill
\begin{subfigure}{0.325\linewidth}\raw{scene/sky_noon}{\linewidth}\caption{Noon.}\end{subfigure}\hfill
\begin{subfigure}{0.325\linewidth}\raw{scene/dusk_shift_change}{\linewidth}\caption{Dusk: shift change.}\end{subfigure}
\caption{Phases of the day; the fog colour follows the sky's horizon colour in each.}
\label{fig:phases}
\end{figure}

Fog and horizon agree because both come from the same analytic sky (Preetham model \cite{preetham1999}) evaluated for
the current sun direction.
\begin{table}[htbp]
\caption{Everything that follows the sun's elevation.}
\centering\small
\begin{tabularx}{\linewidth}{@{}lX@{}}
\toprule
Derived quantity & Rule\\\midrule
Key light & sun above $-1^\circ$ ($3.2\,s(e)\,T(e)$, direction quantised to $1^\circ$ for the shadow cache); else the moon, fading in as $1-s_{[-6,-1]}(e)$\\
Sky, ambient, fog & by day: a Preetham sky for the sun direction; its spherical-harmonic irradiance is the ambient light; fog takes the horizon colour; night presets blended in log space\\
Exposure & key value moves from the night to the day value as $s_{[-8,10]}(e)$\\
Practical lamps, searchlights & lamps on as $1-s_{[-4,2]}(e)$; UAV searchlights only below $-2^\circ$\\
Rules & intruder waves only at night (unless day ops); soldiers' day shift 06:24--17:36\\
\bottomrule
\end{tabularx}
\end{table}
\textbf{Transitions.} The clock runs automatically (one day = \daySeconds\,s) or is scrubbed by hand (\key{,} \key{.}),
jumps to the next phase (\key{/}) or to a preset in the pause menu; lighting follows continuously in every case.

\textbf{Limitation.} The moon is not an orbiting body: it is a fixed-direction night key light (intensity \moonI) and
does not rise as the sun sets.
